\scanpage{9}
% L1–7
C'est un ens qui a $\unsure{6|\sigma}$ éléments, et on a une \deleted{bijection} surjection
\[
\begin{aligned}
S'&\longrightarrow\unsure{\tau'|T'}\overset{\text{déf}}{=}\bigwedge_i\omega_{S_i}\\
&\hspace{4em}\text{(noté $\omega$ précédemment)}
\end{aligned}
\]
\deleted{\unsure{d}}qui est de degré $3$. Donc $S'$ est \unsure{égalemt|également} un ens. muni d'une partition de type $(3,3)$.

% L8–12
La donnée d'une structure \unsure{hexagonale|hexagonal} équivaut à celle d'une structure $(S,R)$, et d'un $s\in S'$ -- celle d'un hexagone et d'un sommet dessus, à un
% archive sheet number: 8

